\chapter{Formal geometry and cyclic descent}
\label{mht:ch:cyclic-descent}

A local formula for an observable must transform when the Darboux
chart changes. A local comparison must transform with it. The two
transformations can agree up to homotopy on each overlap and still
require another homotopy on triple overlaps. Quantization adds the
period of the algebra and the normalization of its trace to these
gluing data.

The circle comparison of Chapter~\ref{mht:ch:integral-cyclic}
provides a local map and two Hamiltonian homotopies. Its logarithmic
torus example also gives a global comparison on a specified algebra.
A general surface requires coordinate changes and allowed overlap
corrections. Their successive choices need a topology of convergence. The first obstruction
already appears for two one-dimensional complexes.

\section{Coordinates that preserve the formal point}

Put $R=\mathbb C[[q,p]]$, with maximal ideal $\mathfrak m=(q,p)$
and its $\mathfrak m$-adic topology. A pointed formal coordinate
change is a continuous $\mathbb C$-algebra automorphism of $R$.
It preserves $\mathfrak m$. It is symplectic when its pullback
preserves $\omega=dq\wedge dp$.
For maps $\phi$ followed by $\psi$, pullbacks satisfy
\[
 (\psi\circ\phi)^*=\phi^*\psi^*.
\]
Thus the order of pullback operators is fixed by their action on
functions, independently of the notation for a transition matrix.

The Hamiltonian convention remains
$X_f=\{f,-\}$, with $\{q,p\}=1$. It gives
\[
 X_f=f_q\partial_p-f_p\partial_q,
 \qquad -\iota_{X_f}\omega=df.
\]
A symplectic formal vector field $v$ therefore determines a closed
formal one-form $-\iota_v\omega$.

\begin{proposition}\label{mht:prop:formal-coordinate-hamiltonians}
Every symplectic formal vector field is Hamiltonian, uniquely modulo
constants. It vanishes at the formal point exactly when its
Hamiltonian has no linear term. A Hamiltonian in $\mathfrak m^3$
has a convergent formal exponential $e^{X_f}$ acting on $R$.
A nonzero constant translation is not a pointed formal
automorphism.
\end{proposition}

\begin{proof}
Let $E=q\partial_q+p\partial_p$. Decompose a closed formal one-form
as $\alpha=\sum_{r\geq1}\alpha_r$, where pullback by scalar dilation
multiplies $\alpha_r$ by its $r$th power. Cartan's identity gives
\[
 d(\iota_E\alpha_r)=r\alpha_r.
\]
Hence $f=\sum_{r\geq1}r^{-1}\iota_E\alpha_r$ satisfies $df=\alpha$.
The sum converges in the formal topology.

A function with zero
differential is constant.
The constant coefficients of $v$ are
exactly the coefficients of the linear part of $f$.

For $f\in\mathfrak m^3$, the derivation $X_f$ sends
$\mathfrak m^s$ into $\mathfrak m^{s+1}$. Each finite jet of
$e^{X_f}a$ therefore uses finitely many terms. The derivation rule
makes this exponential multiplicative, and $e^{-X_f}$ is its
inverse. It preserves the Poisson bracket. In contrast,
$p\mapsto p+c$ with $c\ne0$ does not preserve the maximal ideal,
so it cannot define a pointed formal automorphism.
\end{proof}

Translations do act after adjoining a formal motion parameter.
For example $e^{t\partial_p}$ acts continuously on
$\mathbb C[[q,p,t]]$ with its $(q,p,t)$-adic topology.
Evaluating $t$ at a nonzero complex number is an additional
operation, unavailable on an arbitrary formal series. Quadratic
Hamiltonians instead generate the linear symplectic group. Their
exponentiation is distinct from the pronilpotent exponential of
$\mathfrak m^3$.

Now let $(X,\omega)$ be a holomorphic symplectic surface and let $v$
be a global holomorphic symplectic vector field. Define
$\alpha_v=-\iota_v\omega$. On sufficiently small coordinate balls
choose holomorphic functions $h_i$ satisfying $dh_i=\alpha_v$.
On each connected component of an overlap, $h_j-h_i$ is constant.
These differences define
\begin{equation}\label{mht:eq:hamiltonian-period}
 \operatorname{per}_X(v)=[h_j-h_i]\in H^1(X,\mathbb C_X).
\end{equation}

Equivalently this is the connecting class for the exact sequence
of sheaves
\[
 0\longrightarrow\mathbb C_X\longrightarrow\mathcal O_X
 \xrightarrow{d}\Omega^1_{X,\mathrm{cl}}\longrightarrow0.
\]
Here $\Omega^1_{X,\mathrm{cl}}$ is the sheaf of closed holomorphic
one-forms. Local exactness follows from the holomorphic Poincar\'e
lemma on a ball. An arbitrary Stein open need not have this global
exactness.

\begin{proposition}\label{mht:prop:global-hamiltonian-period}
The field $v$ has a global holomorphic Hamiltonian if and only if
$\operatorname{per}_X(v)=0$. If $X$ is compact and connected, this
holds if and only if $v=0$.
\end{proposition}

\begin{proof}
Vanishing of the connecting class means that constants can be
added to the $h_i$ so that they agree. The resulting global function
is a Hamiltonian for $v$. A global Hamiltonian conversely supplies
such a choice. On a compact connected complex manifold every
holomorphic function is constant, so its Hamiltonian vector field
is zero. Nondegeneracy of $\omega$ proves the remaining implication.
\end{proof}

For the logarithmic torus, $\omega=d\log x\wedge d\log y$ and
$v=x\partial_x$ give $\alpha_v=-d\log y$. Its integral around the
second circle is $-2\pi i$, so no single-valued Hamiltonian exists.
The vector field and the global cyclic comparison nevertheless
exist. The period obstructs a Hamiltonian function with the stated
coefficients. Global algebras and non-Hamiltonian vector fields can
still exist on the surface.

\section{An obstruction relative to fixed local maps}

Take the complex $\mathbb C$ in degree zero on two charts and their
overlap, with identity restrictions. Prescribe comparison maps
given by multiplication by $1$ and by $2$. They are separately
isomorphisms. Their overlap discrepancy is the element $1$ of
\[
 T^0=\mathbb C^2,\qquad T^1=\mathbb C,
 \qquad D(a,b)=b-a.
\]

It is a boundary in $T$. Removing it while keeping both prescribed
chart maps requires a correction with zero chart components, and
there is no such degree-zero correction. The coefficient space of
allowed corrections is therefore part of the obstruction problem.

Let $\pi:(T,D)\to(Q,d_Q)$ be a degreewise split surjection of
cochain complexes over a commutative ring. The quotient $Q$ records
the local data that must be preserved. Put $K=\ker\pi$, and choose
a graded section $s:Q\to T$. For topological complexes all these
maps, including the section, are required to be continuous.

Define the degree-one map
\[
 b=D s-s d_Q:Q\longrightarrow K[1].
\]
The notation indicates its degree. As operators,
$D b+b d_Q=0$. For a prescribed cycle $\theta\in Q^0$, the defect
of its chosen lift is
\[
 \Omega=D(s\theta)=b\theta\in K^1.
\]

\begin{proposition}\label{mht:prop:relative-descent}
The class $[\Omega]\in H^1(K)$ is independent of the section.
There is a closed lift of $\theta$ to $T^0$ exactly when this class
vanishes. If lifts exist, their classes modulo homotopies in
$K^{-1}$ form a torsor for $H^0(K)$.
\end{proposition}

\begin{proof}
The identity $D^2=0$ makes $\Omega$ closed. A second section changes
$s\theta$ by an element of $K^0$, so it changes $\Omega$ by a
boundary in $K$. Every lift is $s\theta+h$, with $h\in K^0$.
It is closed precisely when $Dh=-\Omega$.

The difference of two closed lifts is a degree-zero cycle in $K$.
Two such differences represent the same relative homotopy class
precisely when they differ by $D(K^{-1})$. This gives the free
transitive action of $H^0(K)$.
\end{proof}

In the two-chart example, $Q=\mathbb C^2$ lies in degree zero,
$K^0=0$, and $K^1=\mathbb C$. Thus $[\Omega]=[1]$ is nonzero in
$H^1(K)$ although $H^1(T)=0$. The absolute complex loses the
requirement that the two local maps remain fixed.
Existence need not make the choices contractible. For
$K=\mathbb C$ in degree zero, distinct elements give distinct
homotopy classes of lifts.

Sometimes local changes are permitted. Specify them by a complex
$L$ and a chain map $j:L\to Q$. An allowed change replaces
$\theta$ by $\theta+jw$, with $w\in Z^0(L)$.

Put
\begin{equation}\label{mht:eq:allowed-change-complex}
 E^n=K^n\oplus L^n,
 \qquad d_E(k,w)=(Dk+bjw,d_Lw).
\end{equation}
The relation $Db+bd_Q=0$ makes $d_E^2=0$.
The class $[(\Omega,0)]\in H^1(E)$ vanishes exactly when
\[
 d_Lw=0,\qquad Dh+bjw=-\Omega
\]
has a solution. These equations say that $s(\theta+jw)+h$ is a
closed lift.
On the two charts, diagonal changes $(w,w)$ leave the
discrepancy unchanged, whereas allowing arbitrary changes in
$\mathbb C^2$ removes it. The two choices define different
comparison problems.

\section{Overlap equations and their higher terms}

For cochain complexes $P,G$, the differential on a homogeneous map
$f:P\to G$ of degree $r$ is
\[
 d_{\mathrm{Hom}}f=d_Gf-(-1)^r f d_P.
\]
Suppose first that $P,G$ are complexes of sheaves on an ordered
cover $\{U_i\}$. Their continuous Hom sheaf $\mathcal H$ has the
product \v Cech total complex
\begin{equation}\label{mht:eq:cech-totalization}
 T^n=\prod_{p\geq0}\check C^p(\mathfrak U,\mathcal H^{n-p}),
 \qquad D=\check\delta+(-1)^p d_{\mathrm{Hom}}.
\end{equation}
Projection to \v Cech degree zero is a chain map onto the complex
of chart data. Its kernel is the relative complex of
Proposition~\ref{mht:prop:relative-descent}.

A closed element of total degree zero has components $\Phi_p$ of
internal degree $-p$. Its equations begin
\[
 d\Phi_0=0,\qquad d\Phi_1=\check\delta\Phi_0,
 \qquad d\Phi_2=-\check\delta\Phi_1.
\]
Each equation is a primitive problem in its actual Hom complex.
To identify this totalization with derived global sections requires
an acyclic resolution or the corresponding hypotheses on the cover
and totalization. The displayed algebraic equations hold without
that further identification.

Object transitions may themselves be coherent only up to homotopy.
Write $u_{ij}:P_i\to P_j$ and $v_{ij}:G_i\to G_j$ for chain maps
on overlaps. All compositions below are restricted to the full
intersection. Choose degree-minus-one object homotopies satisfying
\[
 dt^P_{ijk}=u_{jk}u_{ij}-u_{ik},\qquad
 dt^G_{ijk}=v_{jk}v_{ij}-v_{ik}.
\]

For local comparison maps $\Theta_i:P_i\to G_i$, a pair homotopy
$h_{ij}:P_i\to G_j$ must satisfy
\begin{equation}\label{mht:eq:descent-pair}
 dh_{ij}=\Theta_j u_{ij}-v_{ij}\Theta_i.
\end{equation}
The right side is closed. Its exactness is an additional condition.

\begin{lemma}\label{mht:lem:descent-triple}
If the pair equations hold, then
\begin{equation}\label{mht:eq:descent-triple}
 c_{ijk}=v_{jk}h_{ij}+h_{jk}u_{ij}-h_{ik}
              +t^G_{ijk}\Theta_i-\Theta_k t^P_{ijk}
\end{equation}
is a degree-minus-one cycle. The next comparison homotopy has degree
$-2$ and obeys $dh_{ijk}=-c_{ijk}$ in the totalization convention
\eqref{mht:eq:cech-totalization}.
\end{lemma}

\begin{proof}
Applying $d_{\mathrm{Hom}}$ to the first three terms and using
\eqref{mht:eq:descent-pair} gives
\[
 \Theta_k(u_{jk}u_{ij}-u_{ik})
       -(v_{jk}v_{ij}-v_{ik})\Theta_i.
\]
The differential of the last two terms is its negative. The sign
of the next primitive is the positive internal differential in
\v Cech degree two.
\end{proof}

The object data require their own next equation. For instance
\[
 u_{kl}t^P_{ijk}+t^P_{ikl}-t^P_{jkl}u_{ij}-t^P_{ijl}
\]
is closed: substitution of the four equations for $dt^P$ cancels
each composite. A degree-minus-two object homotopy must kill this
cycle before the next comparison equation is formulated.

For strict sheaves, a total comparison $\Phi$ acts on total
\v Cech cochains by
\[
 (\Phi x)_{i_0\ldots i_n}
 =\sum_{p=0}^n(-1)^{p(n-p)}
   \Phi_{i_0\ldots i_p}(x_{i_p\ldots i_n}).
\]
The sign moves internal degree $-p$ past \v Cech degree $n-p$.
The cup-product Leibniz identity then gives $D\Phi=\Phi D$.
Each output component uses a finite sum, so this formula is
continuous in the product topology when its component maps are.

An algebra comparison also has multiplication inputs. Its
multiplicativity equation is a map on $P\widehat\otimes P$,
not on $P$ alone. Supported collision operations, normal
derivatives and both vacuum insertions likewise contribute their
own inputs and relations.

A deformation complex for all these
operations requires a resolution of the entire operation diagram.
At a linear extension stage, restriction to fixed local operations
must supply the map $T\to Q$. Its kernel specifies the allowed
corrections. The unary Hom calculation does not construct this
larger complex or its supported maps.

\section{Nonlinear corrections and complete filtrations}

A differential graded Lie algebra $\mathfrak g$ has a degree-one
differential $d$ and a degree-zero graded Lie bracket. For
$a\in\mathfrak g^1$, define
\[
 F(a)=da+\tfrac12[a,a],\qquad d_a=d+[a,-].
\]
The Jacobi and derivation identities give
\begin{equation}\label{mht:eq:curvature-bianchi}
 d_a^2x=[F(a),x],\qquad d_aF(a)=0.
\end{equation}
Thus $d_a$ is a differential at a Maurer--Cartan element, meaning
an element with $F(a)=0$. For a general candidate $a$, a correction
$b\in\mathfrak g^1$ must solve
\begin{equation}\label{mht:eq:full-curvature-correction}
 F(a+b)=F(a)+d_ab+\tfrac12[b,b]=0.
\end{equation}

The quadratic term can be decisive. Take degree-one elements $a,b$,
central degree-two elements $c,z$, and set
\[
 da=c,\qquad db=-c,\qquad [b,b]=2z,
\]
with all other differentials and brackets zero. These formulas
define a dg Lie algebra. Then
\[
 F(a)+d_ab=0,\qquad F(a+b)=z\ne0.
\]
A primitive of the linear defect has solved only its linear
equation.

For higher operations use the cohomological $L_\infty$ convention
$|\ell_n|=2-n$. The curvature and its unary Taylor operator are
\[
 F(a)=\sum_{n\geq1}\frac1{n!}\ell_n(a^{\otimes n}),\qquad
 d_a x=\sum_{n\geq0}\frac1{n!}\ell_{n+1}(a^{\otimes n},x).
\]
These series are finite in a nilpotent algebra, or converge in the
complete filtrations used below.

\begin{proposition}\label{mht:prop:square-zero-descent}
Let $\widetilde{\mathfrak g}\to\mathfrak g$ be a surjective strict
morphism of nilpotent $L_\infty$ algebras over a field of
characteristic zero. Let its kernel be $I$, and assume that every
bracket containing two entries from $I$ vanishes. For a
Maurer--Cartan element $a$ of $\mathfrak g$, choose a graded lift
$\widetilde a$. Then $(I,d_a)$ is a complex and
$F(\widetilde a)\in I^2$ is a closed element. Its class in
$H^2(I,d_a)$ is independent of the lift and vanishes exactly when
a Maurer--Cartan lift exists. Relative gauge classes of such lifts,
when nonempty, form a torsor for $H^1(I,d_a)$.
\end{proposition}

\begin{proof}
The curvature lies in $I$ because its image is $F(a)=0$.
The $L_\infty$ identities give its Bianchi identity. The same
identities express $d_a^2x$ as a sum with both a curvature input
and the input $x\in I$. Every such term vanishes.

Replacing
$\widetilde a$ by $\widetilde a+b$, with $b\in I^1$, leaves $d_a$
unchanged on $I$ and changes the curvature by $d_ab$.
This proves
the obstruction criterion. Differences of lifts are cycles in
$I^1$. Relative degree-zero gauge changes are the boundaries
$d_a(I^0)$. Higher terms in the gauge action contain two kernel
inputs and vanish.
\end{proof}

The obstruction degree is two for these objects and one for the
linear comparison maps of Proposition~\ref{mht:prop:relative-descent}.
A shift changes the degrees of both the elements and the brackets.

Let $\mathfrak g=F^1\mathfrak g$ be complete and separated for a
decreasing filtration. Assume that $d$ preserves it and that
\[
 \ell_k(F^{r_1},\ldots,F^{r_k})
       \subseteq F^{r_1+\cdots+r_k}.
\]
Completeness means $\mathfrak g\simeq\varprojlim_r\mathfrak g/F^r$.
Separatedness means $\bigcap_rF^r=0$.

\begin{theorem}\label{mht:thm:filtered-descent-corrections}
Suppose $a\in F^1\mathfrak g^1$ satisfies $F(a)\in F^r$, with
$r\geq1$. Its leading curvature $o_r$ is a closed degree-two
element of $(\operatorname{gr}^r_F\mathfrak g,d)$.
A correction $b_r\in F^r\mathfrak g^1$ makes the curvature lie in
$F^{r+1}$ exactly when $[o_r]=0$. Successive compatible corrections
at every order converge to a Maurer--Cartan element.
\end{theorem}

\begin{proof}
Each term in the Bianchi identity involving $a$ raises the
filtration, so its leading part is $do_r=0$. The leading correction
equation is
\[
 o_r+d\bar b_r=0.
\]
Every term with two corrections lies in $F^{2r}\subseteq F^{r+1}$,
and every additional background input also raises the filtration.
This proves the equivalence.

The successive corrections form a
Cauchy sequence. Continuity of each bracket identifies its limit
curvature modulo every $F^r$ with zero. Separatedness then makes
the full curvature zero.
\end{proof}

The later classes depend on the earlier choices. Finding unrelated
solutions at every quotient is weaker than constructing this
compatible sequence. Prescribed local operations are retained by
performing the correction in their relative kernel.

For a comparison of coalgebra models with square-zero coderivations
$Q_P,Q_G$, the defect of a degree-zero map $\mathcal X$ is
\[
 \mathcal R=Q_G\mathcal X-\mathcal XQ_P,
 \qquad Q_G\mathcal R+\mathcal RQ_P=0.
\]
At a filtration stage where the new Taylor coefficients enter
linearly, this is a degree-one cycle in the comparison complex.
Changing those coefficients by $h$ changes the defect by $D h$.
The obstruction is its relative $H^1$ class. Nullary terms, unit
conditions and the continuity of every Taylor component remain
part of the chosen model.

\section{Compatible primitives in an inverse system}
\label{mht:sec:descent-towers}

Let $(C_n,d_n)$ be a countable inverse system of cochain complexes
of vector spaces, with chain maps $r_n:C_{n+1}\to C_n$. Products
in this section are algebraic products. A topological application
requires an exact structure with exact products and a Roos
resolution computing its derived inverse limit. All indicated maps
and lifts must belong to its continuous coefficient class.

Write
\[
 Z^q_n=\ker(d_n:C_n^q\to C_n^{q+1}),\qquad
 B^q_n=d_n(C_n^{q-1}),\qquad H^q_n=Z^q_n/B^q_n.
\]
For any tower $A_n$, its Roos operator is
\[
 \Delta_A:\prod_n A_n\longrightarrow\prod_n A_n,
 \qquad \Delta_A(a)_n=a_n-r_na_{n+1}.
\]
Its kernel is $\varprojlim A_n$ and its cokernel is the first
derived inverse limit $\varprojlim^1 A_n$ for this countable
system.

Suppose $o_n\in Z^1_n$ is compatible: $r_no_{n+1}=o_n$.
A strictly compatible primitive is a family $c_n\in C_n^0$ with
\begin{equation}\label{mht:eq:strict-primitive-equations}
 d_nc_n=-o_n,\qquad \Delta_C(c)=0.
\end{equation}
The first equation alone is a separate problem at each level.

\begin{theorem}\label{mht:thm:strict-primitive-tower}
Assume $[o_n]=0$ for every $n$, and choose primitives $c_n$.
Their discrepancies $\delta_n=c_n-r_nc_{n+1}$ lie in $Z^0_n$.
The class
\begin{equation}\label{mht:eq:strict-primitive-class}
 \lambda_Z(o)=[(\delta_n)_n]\in\varprojlim{}^1 Z^0_n
\end{equation}
is independent of the chosen primitives. A strict primitive family
exists if and only if $\lambda_Z(o)=0$.
Its image $\lambda_H(o)$ in $\varprojlim^1 H^0_n$ is sufficient
when $\varprojlim^1 B^0_n=0$.
\end{theorem}

\begin{proof}
Compatibility of $o$ gives $d_n\delta_n=0$.
Changing $c_n$ by a cycle $a_n$ changes $\delta$ by
$\Delta_Z(a)$. This proves independence. If
$\delta=\Delta_Z(a)$, then $c-a$ satisfies both equations
\eqref{mht:eq:strict-primitive-equations}. The converse follows
by comparing any initial primitives with a strict family.

Apply the kernel--cokernel exact sequence to the two product rows
of $0\to B^0_n\to Z^0_n\to H^0_n\to0$ and their Roos operators.
It contains
\[
 \varprojlim{}^1 B^0_n\longrightarrow
 \varprojlim{}^1 Z^0_n\longrightarrow
 \varprojlim{}^1 H^0_n\longrightarrow0.
\]
The first vanishing hypothesis therefore makes the last map an
isomorphism, proving the final assertion.
\end{proof}

Thus a cohomology-level discrepancy must not be substituted for a
cycle-level discrepancy without a condition allowing its boundary
remainders to be removed. This distinction disappears for complexes
concentrated in nonnegative degrees, since then $B^0_n=0$.

A tower is Mittag--Leffler if, for each fixed $n$, the images of
$A_m\to A_n$ stabilize as $m$ increases. A countable
Mittag--Leffler tower has surjective Roos operator.

To see this, take its stable-image tower. Its transition maps are
surjective, so $a_n-r_na_{n+1}=b_n$ can be solved recursively.
The quotient tower has eventually zero maps to each fixed level.
There the solution is the finite sum
\[
 a_n=b_n+r_nb_{n+1}+r_nr_{n+1}b_{n+2}+\cdots.
\]
Lifting the quotient solution and correcting its discrepancy in
the stable-image tower solves the original equation.

Finite-dimensional towers are Mittag--Leffler. A finite cutoff on
field labels does not by itself make a completed functional space
finite-dimensional.

\subsection*{The derived inverse limit}

The two-column Roos complex of the tower $C_n$ has terms and
differential
\begin{equation}\label{mht:eq:derived-roos-complex}
 \begin{split}
 \mathcal R(C)^k&=\prod_n C_n^k\oplus\prod_n C_n^{k-1},\\
 d_{\mathcal R}(x,y)&=(dx,\Delta_Cx-dy).
 \end{split}
\end{equation}
The chain-map identity $d\Delta_C=\Delta_Cd$ proves
$d_{\mathcal R}^2=0$. It is the shifted cone of the Roos operator.
Its long exact cohomology sequence gives
\begin{equation}\label{mht:eq:milnor-roos-sequence}
 0\longrightarrow\varprojlim{}^1 H^{k-1}_n
 \longrightarrow H^k(\mathcal R(C))
 \longrightarrow\varprojlim H^k_n\longrightarrow0.
\end{equation}
Products of vector spaces are exact, so cohomology commutes with
the products used in this derivation. No splitting of this exact
sequence is specified.

The compatible residual gives the closed element $(o,0)$ in degree
one. A primitive in $\mathcal R(C)$ consists of $c_n\in C_n^0$
and $h_n\in C_n^{-1}$ satisfying
\[
 dc=-o,\qquad \Delta_Cc=dh.
\]
These are homotopy-compatible primitives. After solving the first equation, alter $c$ by cycles.
The second equation is then solvable exactly when its discrepancy
vanishes in $H^0$, which is the condition $\lambda_H(o)=0$.

The class in the left term of
\eqref{mht:eq:milnor-roos-sequence} is therefore secondary to the
vanishing of the image in its right term. The exact sequence does
not define a canonical pair of obstruction coordinates beforehand.

If each degreewise Roos operator $\Delta_{C^k}$ is surjective,
the inclusion $\varprojlim C_n\to\mathcal R(C)$ is a
quasi-isomorphism. Indeed the short exact sequence of complexes
\[
 0\longrightarrow\varprojlim C_n\longrightarrow
 \prod_n C_n\xrightarrow{\Delta_C}\prod_n C_n\longrightarrow0
\]
identifies its kernel with the homotopy fiber.
This supplies another sufficient condition for replacing a
homotopy-compatible primitive by a strict one.

\begin{example}\label{mht:ex:strict-versus-derived-primitives}
For $n\geq1$, consider the acyclic complex
\[
 C_n^{-1}=x^n\mathbb C[x]\ \lhook\joinrel\longrightarrow\
 C_n^0=\mathbb C[x]\ \longrightarrow\
 C_n^1=\mathbb C[x]/(x^n),
\]
with inclusion and quotient as differentials. Transitions are
inclusion in degree $-1$, identity in degree zero, and quotient in
degree one. The compatible residual
\[
 o_n=1+x+\cdots+x^{n-1}\pmod{x^n}
\]
has primitives $c_n=-(1+x+\cdots+x^{n-1})$. Their discrepancies are
$c_n-c_{n+1}=x^n\in Z^0_n=x^n\mathbb C[x]$.

A strict primitive would be one polynomial with every Taylor
coefficient equal to $-1$. It would satisfy $(1-x)c=-1$.
Evaluation at $x=1$ proves that no such polynomial exists.

Nevertheless $H^0_n=0$ at every level, and
$h_n=x^n\in C_n^{-1}$ satisfies $dh_n=x^n$.
Thus $(c,h)$ is a primitive in the Roos complex.
Replacing the degree-zero polynomial ring by $\mathbb C[[x]]$
admits $c=-(1-x)^{-1}$ and changes the strict problem.
The example distinguishes both kinds of compatibility and the
coefficient completion that relates them.
\end{example}

\subsection*{Compatible representatives and closed exchange}

Suppose a compatible cohomology class
$\alpha\in\varprojlim H^1(X_n)$ has representatives
$w_n\in Z^1(X_n)$. Their discrepancies
\[
 w_n-r_nw_{n+1}\in B^1(X_n)
\]
define a class
\begin{equation}\label{mht:eq:strict-representative-class}
 \mu_B(\alpha)\in\varprojlim{}^1 B^1(X_n).
\end{equation}
Changing a representative adds a boundary, and hence changes its
discrepancy by a Roos coboundary in the boundary tower. It follows
that $\mu_B$ is well-defined and vanishes precisely when
$\alpha$ has strict compatible cycle representatives.

Writing $w_n-r_nw_{n+1}=du_n$ does not put $u_n$ in $Z^0(X_n)$.
In $\mathbb C\xrightarrow{1}\mathbb C$, placed in degrees zero
and one, the degree-one cycles $1$ and $0$ differ by $d(1)$.
Their degree-zero primitive is not closed.
There is no class $[u_n]\in H^0(X_n)$ to use in a Roos cokernel.

This gives an explicit strict criterion for a closed-field
contribution together with a counterterm. Let
$\Xi_n:X_n\to Q_n$ be compatible chain maps and let
$o_n\in Z^1(Q_n)$ be compatible. The equations are
\begin{equation}\label{mht:eq:closed-exchange-tower}
 o_n+\Xi_n W_n+dC_n=0,\qquad dW_n=0,
\end{equation}
with $W_n$ and $C_n$ both strictly compatible.
Define the cone
\[
 E_n^k=Q_n^k\oplus X_n^{k+1},\qquad
 d_E(q,x)=(d_Qq+\Xi_nx,-d_Xx).
\]
Then \eqref{mht:eq:closed-exchange-tower} asks for a strict
primitive of $(o_n,0)$ in $E_n$. Theorem
\ref{mht:thm:strict-primitive-tower} applies directly.

The same criterion can be resolved into successive choices.
First choose
\[
 \alpha\in\varprojlim H^1(X_n),\qquad
 H^1(\Xi_n)(\alpha_n)=-[o_n].
\]
Next require $\mu_B(\alpha)=0$ and choose compatible representatives
$W_n$. Finally require
$\lambda_Z(o+\Xi W)=0$ in $\varprojlim^1 Z^0(Q_n)$.

These conditions are necessary and sufficient by their definitions
and the two primitive criteria. The last condition can depend on
the preceding choice of $W$. If the towers $B^1(X_n)$ and
$Z^0(Q_n)$ are Mittag--Leffler, only the first choice remains.
Neither of these cycle conditions is an automatic consequence of
Mittag--Leffler cohomology groups.

\section{Linear transport at an interacting background}

Renormalization acts on interactions nonlinearly. To transport an
obstruction complex one needs its derivative at the chosen
interaction. The distinction is already visible in a finite
Gaussian model.

Let $P$ be a constant symmetric covariance on a finite-dimensional
even field space and put
$\partial_P=\frac12P^{ij}\partial_i\partial_j$.
Use polynomial interactions with a positive formal coupling, so
that all expressions are interpreted coefficientwise. Define
\[
 Z_P(I)=e^{\hbar\partial_P}e^{I/\hbar},\qquad
 W_P(I)=\hbar\log Z_P(I).
\]

The exponential and logarithm here specify the connected formal
graph expansion. At a fixed coupling order there are finitely many
vertices. Polynomial valences make each graph sum finite. In the
connected expansion every power of $\hbar$ is a loop number and
is nonnegative.

\begin{proposition}\label{mht:prop:rg-tangent-transport}
The derivative of $W_P$ at $I$ is the linear operator
\begin{equation}\label{mht:eq:rg-tangent}
 T_{P,I}(F)=Z_P(I)^{-1}e^{\hbar\partial_P}
                           (F e^{I/\hbar}).
\end{equation}
For another covariance $Q$,
\[
 T_{P+Q,I}=T_{Q,W_P(I)}\,T_{P,I}.
\]
If the Gaussian transport intertwines the master-equation vector
fields and $I$ solves that equation, $T_{P,I}$ is a chain map
between their tangent complexes.
\end{proposition}

\begin{proof}
Differentiate $\hbar\log Z_P(I+sF)$ at $s=0$ to obtain
\eqref{mht:eq:rg-tangent}. Constant contraction operators commute,
so
$Z_Q(W_P(I))=e^{\hbar\partial_Q}Z_P(I)=Z_{P+Q}(I)$.
Differentiating this identity proves the composition formula.
For the last assertion, differentiate the stated intertwining
identity at a zero of the master-equation vector field. Its
derivatives there are the tangent differentials, so the resulting
identity is precisely the chain-map equation.
\end{proof}

For one coordinate, covariance $P$, and $I=gx^2/2$, the formal
Gaussian identity gives
\[
 W_P(gx^2/2)
 =\frac{gx^2}{2(1-Pg)}-\frac{\hbar}{2}\log(1-Pg).
\]
It follows either by completing the square formally or by solving
the heat equation with the given initial exponential.
The term $Pg^2x^2/2$ shows that $W_P$ is not linear in $I$.
Consequently $W_P(C)$ is not the transition map on a linear
counterterm complex.

\begin{example}[RG can produce nonlocal terms]
\label{mht:ex:rg-nonlocal}
Consider two sites, with fields $x,y$ and local functionals
$\C[x]+\C[y]\subset\C[x,y]$. Choose
$\partial_P=\partial_x\partial_y$ and the local interaction
$I=g(x^3+y^3)/6$. Direct expansion gives
\[
 W_P(I)=\frac g6(x^3+y^3)
   +g^2\left(\frac{x^2y^2}{4}
             +\frac{\hbar xy}{2}+\frac{\hbar^2}{6}\right)
   \pmod{g^3}.
\]

Indeed, the mixed part of the coefficient of $g^2$ in
$e^{I/\hbar}$ is $x^3y^3/(36\hbar^2)$. Its successive cross
derivatives have coefficients $9,36,36$. In the heat exponential
these are multiplied by $\hbar$, $\hbar^2/2$ and $\hbar^3/6$.

The logarithm cancels the uncontracted quadratic term, giving the
displayed expression after multiplication by $\hbar$.

The term $g^2x^2y^2/4$ depends on both sites and lies outside the
local subspace. Thus finite-scale transport requires a space of
effective functionals larger than the local counterterms. The
locality condition for a field theory must specify the allowed
small-scale asymptotics and singular counterterms.
\end{example}

Fix compatible lower-order interactions and a filtration order
$n$. If $C_n\in F^n$, then
\[
 W_P(I_{<n}+C_n)
 =W_P(I_{<n})+T_{P,I_{<n}}C_n\pmod{F^{n+1}}.
\]
This statement requires the filtration to make terms with two new
corrections lie in $F^{n+1}$. The order-$n$ mismatch of the already
chosen lower interactions remains an inhomogeneous term. The
linear transition on the next correction complex is the induced
$T_{P,I_{<n}}$, with that mismatch retained separately.

For local field theories one must construct continuous transport
on the specified effective functionals. Its small-scale locality
conditions and boundary identities require further analytic
arguments. The finite Gaussian formula identifies the map whose
properties must be proved.

\section{A Weyl connection with prescribed central curvature}

Coordinate transport also acts on the quantized coefficient
algebra. A useful global construction begins with more geometric
data than a symplectic form alone. Suppose that $(X,\omega)$ has a
global holomorphic torsion-free symplectic connection $\nabla$.
This is a hypothesis. Local Darboux connections do not themselves
give such a global connection.

Let $\mathcal W$ be the bundle of formal Weyl algebras on $T_X$.
In a symplectic chart use fiber variables $Q,P$, with
$[Q,P]_\star=\epsilon$. Complete by total weight, assigning weight
one to $Q,P$ and weight two to $\epsilon$. Base holomorphic forms
have weight zero. On $\Omega_X^\bullet(\mathcal W)$ define
\[
 \delta=dq\,\partial_Q+dp\,\partial_P.
\]

For a monomial of fiber degree $s$ and form degree $t$, put
\[
 \delta^{-1}a=
 \begin{cases}
  (s+t)^{-1}(Q\iota_{\partial_q}+P\iota_{\partial_p})a,
       &s+t>0,\\
  0,&s=t=0.
 \end{cases}
\]

These operators are coordinate independent. If $\sigma$ projects
to fiber degree and form degree zero, then
\begin{equation}\label{mht:eq:fedosov-contraction}
 \delta\delta^{-1}+\delta^{-1}\delta=1-\sigma,
 \qquad (\delta^{-1})^2=0.
\end{equation}
The identity follows by counting the $s+t$ occurrences of a fiber
variable or a base differential.

The connection extends as a derivation of the Weyl product.
Torsion-freeness and preservation of $\omega$ give
\[
 [\delta,\nabla]=0,\qquad
 \nabla^2=\epsilon^{-1}[R,-]_\star,
 \qquad \delta R=0,\qquad\nabla R=0,
\]
where $R$ is the quadratic fiber representative of its curvature,
a two-form of total weight two. Brackets of form-valued elements
use the graded product, including exterior multiplication.
Division of a Weyl commutator by $\epsilon$ is integral: the
commutator is divisible by $\epsilon$.

Fix a closed central series
$\Omega\in\epsilon\Omega^2_{X,\mathrm{cl}}[[\epsilon]]$.
The word central means that its coefficients contain no fiber
variables.

\begin{theorem}\label{mht:thm:holomorphic-fedosov-resolution}
There is a unique one-form $r$ of weight at least three satisfying
$\delta^{-1}r=0$ and
\begin{equation}\label{mht:eq:fedosov-recursion}
 r=\delta^{-1}\left(R-\Omega+\nabla r
                         +\frac1{2\epsilon}[r,r]_\star\right).
\end{equation}
The derivation
\[
 D=-\delta+\nabla+\epsilon^{-1}[r,-]_\star
\]
is flat. Projection $\sigma$ identifies its sheaf of flat
degree-zero sections with $\mathcal O_X[[\epsilon]]$ as modules.
The Weyl product transports to an associative deformation product
on this module, with first commutator $\epsilon\{f,g\}$.
The complex $(\Omega_X^\bullet(\mathcal W),D)$ has no positive
cohomology sheaves.
\end{theorem}

\begin{proof}
The operator $\delta^{-1}\nabla$ raises weight by one.
For $r$ of weight at least three,
$\delta^{-1}\epsilon^{-1}[r,-]_\star$ raises weight by at least
two. Iteration of \eqref{mht:eq:fedosov-recursion}, starting with
zero, therefore determines each weight after finitely many
steps and gives the unique solution. It has weight at least three
and satisfies the stated normalization by $(\delta^{-1})^2=0$.

Put
\[
 F=R-\delta r+\nabla r+\frac1{2\epsilon}[r,r]_\star.
\]
Expansion gives $D^2=\epsilon^{-1}[F,-]_\star$. The connection
identities give $DF=0$. Since $\Omega$ is central and closed,
$D(F-\Omega)=0$ as well.

The recursion and
\eqref{mht:eq:fedosov-contraction} imply
$\delta^{-1}(F-\Omega)=0$. Its form degree is two, so its
$\sigma$ projection vanishes. It consequently satisfies
\[
 F-\Omega
 =\delta^{-1}\left(\nabla(F-\Omega)
               +\epsilon^{-1}[r,F-\Omega]_\star\right).
\]
The right side raises weight. The only solution in the separated
completion is zero. Thus $F=\Omega$ is central and $D^2=0$.

For $f\in\mathcal O_X[[\epsilon]]$, solve
\begin{equation}\label{mht:eq:fedosov-flat-lift}
 \tau(f)=f+\delta^{-1}\left(\nabla\tau(f)
                         +\epsilon^{-1}[r,\tau(f)]_\star\right).
\end{equation}
The same weight argument gives a unique solution. Its defect
$D\tau(f)$ has zero $\delta^{-1}$ component, is $D$-closed by
flatness, and has positive form degree. The preceding contraction
argument makes it zero. Conversely a flat section with symbol $f$
satisfies \eqref{mht:eq:fedosov-flat-lift}, proving uniqueness.

Define $f\star_X g=\sigma(\tau(f)\star\tau(g))$.
The product of flat sections is flat, since $D$ is a derivation.
Uniqueness of the lift gives associativity. Modulo $\epsilon$ and fiber degree at least two,
$\tau(f)=f+Qf_q+Pf_p$. Their single contraction gives
the first commutator $\epsilon\{f,g\}$.

Finally let $a$ be $D$-closed of positive form degree. Solve
\[
 b=-\delta^{-1}a+
        \delta^{-1}\left(\nabla b+\epsilon^{-1}[r,b]_\star\right),
 \qquad \sigma b=0.
\]
Again the recursion converges. The error $Db-a$ is closed and
has zero $\delta^{-1}$ component and zero symbol, hence is zero
by the same weight argument. This proves positive-degree
acyclicity.
\end{proof}

This is the Weyl version of the Fedosov contraction. Its formal
resolution method also underlies the covariant formality
construction of Dolgushev, \emph{Covariant and equivariant formality
theorems}, \S3. The theorem above uses holomorphic data supplied
globally. A smooth resolution or a Dolbeault resolution has
different coefficients and requires its own comparison to
holomorphic sections.

\subsection*{Period and algebroid gluing}

The central curvature retains information invisible to the flat
derivation $D^2=0$. In a Darboux chart the linear one-form
\[
 \gamma_{-1}=P\,dq-Q\,dp
\]
satisfies
$\epsilon^{-1}[\gamma_{-1},-]_\star=-\delta$ and
$\epsilon^{-1}\gamma_{-1}\star\gamma_{-1}=\omega$.
With this convention the scalar Weyl curvature is
$\omega+\Omega$, and the associated period is
\begin{equation}\label{mht:eq:weyl-period}
 \vartheta=\epsilon^{-1}[\omega+\Omega]
 \in\epsilon^{-1}[\omega]+H^2_{\mathrm{dR}}(X,\mathbb C)[[\epsilon]].
\end{equation}

The class is understood in holomorphic de Rham hypercohomology
when no global form representative is available. A central
one-form added to a connection representative changes its scalar
curvature by an exact form. Thus such a change preserves
\eqref{mht:eq:weyl-period}. A comparison allowing a different
period is a different quantization problem. Flatness alone does
not erase this parameter.

When no global holomorphic connection has been specified, one
can instead ask for local quantized algebras $A_i$ and algebroid
descent data. In a strict presentation these consist of
isomorphisms $g_{ij}:A_j\to A_i$ and invertible elements
$a_{ijk}\in A_i$ with
\begin{align*}
 g_{ij}g_{jk}&=\operatorname{Ad}(a_{ijk})g_{ik},\\
 a_{ijk}a_{ikl}&=g_{ij}(a_{jkl})a_{ijl}.
\end{align*}

All terms are restricted to the indicated intersection.
The first equation identifies the two ways of transporting a
module on a triple overlap. The second makes that identification
coherent on quadruple overlaps. A strict algebra sheaf is the
special case $a_{ijk}=1$. Local existence of a flat Weyl algebra
does not choose these intertwiners or prove their coherence.

Holomorphic quantization and its characteristic classes are
treated by Nest--Tsygan, \emph{Deformations of symplectic Lie
algebroids, deformations of holomorphic symplectic structures,
and index theorems}, \S5. The local connection calculation above
provides a construction on the stated class of surfaces with a
global holomorphic symplectic connection. A general surface
requires the additional descent or resolution data.

\subsection*{The density and its unit value}

A cyclic density is further structure: a map from the chosen cyclic
chain complex to a specified complex of forms or currents.
It must also have coherent overlap homotopies. A closed integration
functional in the required degree, defined on the chosen supports,
then gives the corresponding cyclic functional.

The product and period do not fix its normalization. Multiplying
a trace by $c(\epsilon)\in\mathbb C[[\epsilon]]^\times$ preserves
cyclicity and multiplies its unit value by $c(\epsilon)$.

For the logarithmic torus of
Chapter~\ref{mht:ch:integral-cyclic}, the constant Laurent
coefficient gives the explicit normalization $\tau(1)=1$.
On a general quantization, a unit-index statement must specify
its density map, normalization and supports. It must also fix
the cohomology class in which the unit is evaluated.

The periodic algebraic index theorem relates a localized trace
density to the tangent bundle's characteristic classes and the
quantization period. The relevant result is Theorem~3.3.1 of
Bressler--Nest--Tsygan, \emph{Riemann--Roch theorems via
deformation quantization II}. Its characteristic expression involves the
$\widehat A$ class and the exponential of the period, with the
parameter regrading of that theorem.

It does not by itself give
an integral comparison over $\mathbb C[[\epsilon,u]]$.
The point-current classes in
Chapter~\ref{mht:ch:integral-cyclic} are annihilated by both
parameters. Inverting either parameter loses precisely such
information. Recovering it requires an integral map and cannot
be inferred from an equality after localization.

\section{Descent of the structured comparison}

The preceding constructions separate two questions: whether the
objects glue, and whether the specified comparison glues with all
its operations. State the latter in a category that retains the
operations being compared.

Let $\mathcal C$ be an infinity-category of complete complexes
with a specified algebraic or supported operation structure.
Assume that the required limits exist in $\mathcal C$. Let
$P_\bullet,G_\bullet$ be their descent diagrams on a cover, with
all overlap and higher coherence maps supplied. Effective descent
means that the global objects under consideration are identified
with $\operatorname*{holim}P_\bullet$ and
$\operatorname*{holim}G_\bullet$. A map of diagrams includes
homotopies in every simplicial degree, not merely maps on charts.

\begin{proposition}\label{mht:prop:structured-global-descent}
Under these hypotheses a homotopy-coherent map
$\Phi_\bullet:P_\bullet\to G_\bullet$ induces a global map in
$\mathcal C$. If each component is an equivalence, its global
map is an equivalence. Every operation or marking retained in
$\mathcal C$ is preserved by this map.
\end{proposition}

\begin{proof}
Apply the limit functor to the map of diagrams. A componentwise
equivalence is an equivalence in the diagram infinity-category,
so a functor sends it to an equivalence. Effective descent
identifies the resulting limits with the specified global
objects. Their operations and markings are preserved because
the construction takes place in $\mathcal C$.
\end{proof}

For unary complexes the overlap equations were written explicitly
above. For a cyclic algebra comparison the category must also
retain multiplication, the mixed differential, the coefficient
completion, the unit and the density. For the chiral problem it
must in addition retain the full supported collision diagram and
normal derivatives. The proposition applies once these data and
their descent theorem are supplied. An equivalence of the
underlying complexes does not supply them.

A sheaf of $E_2$ algebras illustrates the variance issue.
Fix a holomorphic open $U$ and an $E_2$ algebra of sections.
Its factorization realization on framed topological two-disks
requires the corresponding theorem in the chosen completed tensor
category.

For $V\subset U$, sheaf restriction goes from sections on $U$ to
sections on $V$. A factorization algebra on the product needs
covariant structure maps for disjoint product opens and their
descent. The sheaf restrictions therefore do not define the
holomorphic part of that factorization algebra.
Extension from
framed disks to oriented disks also requires the appropriate
rotation equivariance. A boundary category and a bulk action
require additional maps even after the bulk realization exists.

The global diagonal bimodule requires a comparison with the local
diagonal resolutions. This comparison is needed to identify the
global Morita center with derived sections of local Hochschild
cochains. The affine derived Morita theorem
of Chapter~\ref{mht:ch:centres-actions} supplies the comparison under its
own generator hypotheses. A general holomorphic algebroid has
no such identification solely from its local Weyl presentation.

The global integral cyclic problem can now be stated with its
data fixed. Choose a holomorphic symplectic surface, a quantized
Weyl algebroid with period $\vartheta$, a normalized cyclic
density, and a coefficient and support completion over
$\mathbb C[[\epsilon,u]]$.

Construct an integral cyclic comparison extending the local circle
map on this algebroid. It must preserve coherent coordinate
transport, the prescribed unit-index class and the mixed theory's
supported operations.

The two-chart
kernel gives its first relative obstruction. Higher overlap
equations and compatible filtered choices remain to be solved
in their respective complexes. Finally a map from the geometric
closed-source complex must identify this cyclic comparison with
the physical one. The logarithmic torus and the global-connection
construction establish specified portions of these data.
The final physical map on a general surface remains to be constructed.
